\documentclass[10pt]{article}
\usepackage[margin=0.76in]{geometry}
\usepackage{amsmath,amssymb,amsthm}
\usepackage[hidelinks]{hyperref}
\setlength{\parindent}{0pt}
\setlength{\parskip}{0.42em}

\newtheorem{theorem}{Theorem}
\newtheorem{corollary}[theorem]{Corollary}
\newtheorem{proposition}[theorem]{Proposition}

\newcommand{\wstar}{\mathrm{w}^*}
\newcommand{\Orb}{\mathcal O}

\begin{document}

\begin{center}
{\LARGE A Compact-Invariant-Core Form of the}\par\vspace{0.15em}
{\LARGE Bounded-Orbit Hypothesis}\par\vspace{0.45em}
{\large A candidate full answer to Remark 3.6 of arXiv:2404.18261}\par\vspace{0.2em}
August 9, 2026
\end{center}

\begin{center}
\fbox{\parbox{0.91\textwidth}{\textbf{Status.}
Candidate full solution, likely valid.  The proof is elementary and the orbit
hull construction is classical in semigroup fixed-point theory; novelty of the
explicit semihypergroup formulation requires expert review.}}
\end{center}

\section*{The question and the answer}

Let $K$ be a semitopological semihypergroup and let $X$ be a Banach space.
Theorem 3.5 of Bandyopadhyay~\cite{Bandyopadhyay2024} characterizes left
amenability of $AP(K)$ by common fixed points of separately continuous,
equicontinuous affine actions
\[
   \pi:K\times (X^*,\wstar)\longrightarrow (X^*,\wstar)
\]
having a norm-bounded orbit.  Remark 3.6 asks for an equivalent, more
conventional replacement for that orbit condition.

There is an exact replacement: require a nonempty weak-* compact convex
invariant subset of $X^*$.  In fact this is equivalent to the original
hypothesis before amenability or equicontinuity enters.

\begin{theorem}[bounded-orbit/invariant-core equivalence]\label{thm:core}
Let $\pi$ be a separately weak-* continuous affine action of a nonempty
semitopological semihypergroup $K$ on $X^*$.  The following conditions are
equivalent.
\begin{enumerate}
\item[(i)] There is $u\in X^*$ for which
  $\Orb(u)=\{\pi(t,u):t\in K\}$ is norm bounded.
\item[(ii)] There is $u\in X^*$ for which $\Orb(u)$ is relatively weak-*
  compact.
\item[(iii)] There is a nonempty weak-* compact convex set $C\subseteq X^*$
  such that $\pi(s,C)\subseteq C$ for every $s\in K$.
\end{enumerate}
When (i) holds, one may take
\[
             C=\overline{\operatorname{co}}^{\wstar}\Orb(u).
\]
\end{theorem}

\begin{proof}
The equivalence of (i) and (ii) is a useful special feature of dual spaces.
If $\Orb(u)$ is norm bounded, it lies in a scalar multiple of the weak-* compact
dual unit ball, so Banach--Alaoglu makes it relatively weak-* compact.
Conversely, the weak-* closure of a relatively weak-* compact orbit is weak-*
compact.  Every evaluation $v\mapsto v(x)$ is therefore bounded on it for each
$x\in X$, and the uniform boundedness principle gives a uniform norm bound.

Assume (i) and set
$C=\overline{\operatorname{co}}^{\wstar}\Orb(u)$.  The orbit is contained in
some $rB_{X^*}$, so $C$ is a weak-* closed subset of the weak-* compact set
$rB_{X^*}$.  Hence $C$ is nonempty, weak-* compact, and convex.

It remains to prove invariance, which is the only semihypergroup-specific
point.  Fix $s,t\in K$.  By the action axiom,
\[
 \pi(s,\pi(t,u))
   =\int_K \pi(\zeta,u)\,d(p_s*p_t)(\zeta).
\]
The convolution $p_s*p_t$ is a compactly supported probability measure.
Since $\zeta\mapsto\pi(\zeta,u)$ is weak-* continuous, its integral is a
weak-* barycenter of points of $\Orb(u)$ and consequently belongs to $C$.
Thus $\pi(s,\Orb(u))\subseteq C$.  Affineness gives
$\pi(s,\operatorname{co}\Orb(u))\subseteq C$, and weak-* continuity of
$\pi_s$ extends this inclusion to $C$.  Therefore (iii) holds.

Finally assume (iii), choose $u\in C$, and use invariance to obtain
$\Orb(u)\subseteq C$.  A weak-* compact subset of $X^*$ is norm bounded by
the uniform boundedness principle, so (i) follows.
\end{proof}

\begin{corollary}[answer to the replacement question]\label{cor:answer}
In Theorem 3.5 of~\cite{Bandyopadhyay2024}, ``admits an element with bounded
orbit'' may be replaced, without changing the class of representations, by
either of the following:
\begin{enumerate}
\item the action admits a relatively weak-* compact orbit;
\item the action admits a nonempty weak-* compact convex invariant subset.
\end{enumerate}
In particular, the second condition is the usual compact invariant-domain
formulation of an affine fixed-point property.
\end{corollary}

\begin{proof}
Apply Theorem~\ref{thm:core} to each representation quantified over in
Theorem 3.5.  Equicontinuity and left amenability are not needed for the
equivalence of the three orbit conditions; they remain exactly as in the
source theorem.
\end{proof}

\section*{The restriction cannot be omitted}

The replacement above is exact, but no orbit condition can simply be erased
from the theorem.

\begin{proposition}[sharpness]\label{prop:sharp}
There is a semitopological semihypergroup $K$ for which $AP(K)$ is left
amenable and a jointly continuous, equicontinuous affine action of $K$ on a
dual Banach space having no common fixed point.
\end{proposition}

\begin{proof}
Take the discrete semihypergroup associated with the additive group
$K=\mathbb Z$, so $p_m*p_n=p_{m+n}$.  The group $\mathbb Z$ is amenable:
weak-* cluster points of its F\o lner averages give a translation-invariant
mean on $\ell^\infty(\mathbb Z)$, whose restriction is a left invariant mean
on $AP(\mathbb Z)$.

Identify $\mathbb R$ with the dual of the one-dimensional real Banach space
and define
\[
                    \pi(n,x)=x+n.
\]
This is a jointly continuous affine action, and its coefficient maps are
isometries, hence form an equicontinuous family.  A common fixed point would
satisfy $x+1=x$, which is impossible.  Each orbit $x+\mathbb Z$ is unbounded,
in agreement with Theorem~\ref{thm:core}.
\end{proof}

\section*{Scope and provenance}

The result fully answers the literal existence question in Remark 3.6: it
gives two equivalent representation conditions, one of which is the standard
compact convex invariant-domain hypothesis.  It does not assert that this
condition is computationally easier in every concrete example.

The use of a weak-* closed convex orbit hull is classical and occurs in the
semigroup precursor of Lau and Zhang~\cite{LauZhang2021}.  The new content
claimed here is only the explicit three-way formulation and the barycentric
verification for the semihypergroup action axiom.  A bounded literature search
found no explicit published statement presenting this as the answer to the
source remark; priority should nevertheless be checked by a specialist.

{\small
\begin{thebibliography}{9}
\bibitem{Bandyopadhyay2024}
C.~Bandyopadhyay,
\emph{Common Fixed Points of Semihypergroup Representations},
arXiv:2404.18261, 2024.

\bibitem{LauZhang2021}
A.~T.-M. Lau and Y.~Zhang,
\emph{Amenability properties and fixed point properties of affine
representations of semigroups},
Proc. Amer. Math. Soc. 149 (2021), 4815--4823.
\end{thebibliography}
}

\end{document}
